\subsection{First Letter}

\begin{quotex}
This is the first of \textbf{Rene Guenon}`s letters to \textbf{Guido de Giorgio}. It shows how things were done prior to Internet blogs and e-mail. The letters were written in French; Guenon was not fluent enough in Italian to write in that language. The criticism of Evola is pretty severe, but Evola was about 27 years old at the time and just transitioning to Guenon style Tradition after several years studying German philosophy. I have left out the salutations and the flowery, but affectionate, closings common at that time
\end{quotex}

Paris, 20 November 1925

I was beginning to get a little nervous since I hadn't received anything from you since your postcard in July, and, fearing that you were still suffering, I thought to write you some words to ask you for some news, when your letter reached me. Since then, I put off writing to you from one day to the next, because I was always very busy since my return here, so that it is I who, in turn, am late answering you. I address this letter to you to Varazze where, I think you told me, you must be settled for some time already.

Thanks for the good idea that you had to invite me to go to meet up with you; however, unfortunately, trips are quite costly at the moment; I hope all the same that we will end up meeting some other day.

What you say about Vulliaud's Kabbalah is quite correct, although perhaps a little severe; basically, it is pretty much what I wrote myself in a more attenuated form. I spoke of that work with different persons who had read it; their assessments agree with ours and are even less enthusiastic. I didn't have a chance to read Vulliaud's new book on the ``Song of Songs", which appeared three or four months ago; it seems to contain fewer discussions and critiques that the Kabbalah, but still too much.

Thanks for the issue of Bilychnis that you sent me; since then, I received from Evola himself a full packet of other reviews containing some of his articles. In acknowledging its reception, I told him that I had several reservations to make on his point of view which seemed to me overly philosophic. He wrote me a rather long and confused letter in which he protested that the philosophical form which he used is for him only a simple means of exposition which does not affect his doctrine itself. I believe none of it, and I persist in thinking that he is really too imbued with philosophy, and especially German philosophy. In an article published by the review Ultra, he alluded to me in a note, in regards to East and West, in terms that prove that he did not understand very much of what I explained; he even went so far as to describe me as a ``rationalist", which is really ridiculous (all the more so since it concerns a book where I expressly asserted the falsity of rationalism!), and that shows clearly that he is one of those who cannot get rid of philosophical labels and who feel the need to apply them wrongly and haphazardly. He told me his intention to write an article on Man and his becoming; I ask myself what that could possibly be; we will certainly see at some point.

Since we are still on Evola, it is still necessary that I tell you that he was hurt by Reghini's criticisms, however, in a very moderate way. He must be rather vain and would like to have nothing but praise; it is true that he is very young. Vulliaud, who doesn't have the same excuse, is almost as susceptible; it appears that he also was rather unhappy with my article; he imagines that only he knows the Kabbalah and is capable of speaking about it. It is to be feared that Evola is soon doing as much for the Tantras, of which he is, however, not very qualified to be dealing with them; he sees all that through his philosophy, from where a type of deformation in the German manner; the true conception of Shakti is a completely different thing from ``voluntarism".

I am happy that you are very interested by the articles of Gnose; the inferiority of the first issues come from the time when I had only a nominal direction; I even had some trouble subsequently clearing out some troublesome people; it would be too long to tell you in detail.

What is that recent book which contains some fragments from the Zohar from Jean de Pauly's translation? I didn't hear anything about it.

I finally made Massignon's acquaintance; he talks a lot, and there is in him a certain affectation, which otherwise is revealed also in his style. It is true that his works are rather difficult to read, and on the other hand, if he has assuredly understood certain things, he has not however penetrated the heart of Muslim esoterism. As for Carra de Vaus, he understands still less, and even from the point of view of exoterism, he ends up committing some gross errors; he is above all a compiler, and the principal utility of his works consists in what in them is brought together of the teachings which are found scattered a little everywhere.

For the Vedanta, up to now, there have not yet been any reviews in the journals; that will always require some time, and furthermore there have been the holidays. Nothing from Masson-Oursel; I don't know if he will speak about it nor what he will say, but I will be quite astonished if it would be better that what he wrote for my earlier book. There were some notes in the journals; I am attaching to this letter copies of the principle ones; you will see that they are not bad. Last month, there was also a long article about me, by Gonzague Truc, in ``Candide"; I currently have a single copy; but I will try to get another one to send you.

Don't delay too long this time in sending back to me news about you.



\flrightit{Posted on 2012-06-18 by Cologero }

\begin{center}* * *\end{center}

\begin{footnotesize}\begin{sffamily}



\texttt{Audrey on 2012-09-17 at 21:15 said: }

You mean there was a time with no internet blogs and e-mail?


\hfill

\texttt{Saladin on 2013-02-14 at 21:22 said: }

I liked the jab at Massignon but was quite uncomfortable reading his views on Evola. I think later on his impression of Evola somehow improved.


\end{sffamily}\end{footnotesize}
